\textbf{Phần 1. Trắc nghiệm lựa chọn}
\setcounterpageref{ex}{0}

\begin{ex} %[3D2V64]
Tính giá trị của phép tính: $84\,650 - 32\,480$.
\choice
{$52\,270$}
{\True $52\,170$}
{$52\,070$}
{$51\,170$}
\loigiai{
$84\,650 - 32\,480 = 52\,170$.
\textbf{Đáp án của câu hỏi là B}
}
\end{ex}

\begin{ex} %[3D2V64]
Một kho gạo có $95\,000\text{ kg}$ gạo. Người ta đã xuất đi $48\,500\text{ kg}$. Trong kho còn lại bao nhiêu ki-lô-gam gạo?
\choice
{$47\,500\text{ kg}$}
{$46\,000\text{ kg}$}
{$56\,500\text{ kg}$}
{\True $46\,500\text{ kg}$}
\loigiai{
$95\,000 - 48\,500 = 46\,500\text{ (kg)}$.
\textbf{Đáp án của câu hỏi là D}
}
\end{ex}

\begin{ex} %[3D2V64]
Tìm số tự nhiên $y$, biết: $80\,000 - y = 35\,600$.
\choice
{\True $44\,400$}
{$45\,400$}
{$44\,600$}
{$54\,400$}
\loigiai{
$y = 80\,000 - 35\,600 = 44\,400$.
\textbf{Đáp án của câu hỏi là A}
}
\end{ex}

\begin{ex} %[3D2V64]
Một đoạn đường sắt dài $85\,400\text{ m}$. Người ta đã nâng cấp được $46\,200\text{ m}$. Đoạn đường sắt còn lại chưa nâng cấp dài bao nhiêu mét?
\choice
{$38\,200\text{ m}$}
{$40\,200\text{ m}$}
{\True $39\,200\text{ m}$}
{$39\,000\text{ m}$}
\loigiai{
Đoạn còn lại: $85\,400 - 46\,200 = 39\,200\text{ (m)}$.
\textbf{Đáp án của câu hỏi là C}
}
\end{ex}

\begin{ex} %[3D2V64]
Tính nhanh giá trị: $74\,500 - 25\,000 - 14\,500$.
\choice
{$40\,000$}
{\True $35\,000$}
{$30\,000$}
{$45\,000$}
\loigiai{
$(74\,500 - 14\,500) - 25\,000 = 60\,000 - 25\,000 = 35\,000$.
\textbf{Đáp án của câu hỏi là B}
}
\end{ex}

\begin{ex} %[3D2V64]
Tìm hiệu của số tròn chục nghìn lớn nhất có 5 chữ số và số tròn chục nghìn bé nhất có 5 chữ số.
\choice
{$90\,000$}
{$70\,000$}
{$100\,000$}
{\True $80\,000$}
\loigiai{
$90\,000 - 10\,000 = 80\,000$.
\textbf{Đáp án của câu hỏi là D}
}
\end{ex}

\begin{ex} %[3D2V64]
Một trang trại có $65\,000$ con vịt. Sau khi xuất bán một số con thì trang trại còn lại $28\,500$ con. Hỏi trang trại đã bán bao nhiêu con vịt?
\choice
{\True $36\,500$ con}
{$37\,500$ con}
{$35\,500$ con}
{$46\,500$ con}
\loigiai{
Số vịt đã bán: $65\,000 - 28\,500 = 36\,500$ con.
\textbf{Đáp án của câu hỏi là A}
}
\end{ex}

\begin{ex} %[3D2V64]
Nếu bớt ở số bị trừ đi 1200 đơn vị và giữ nguyên số trừ thì hiệu sẽ
\choice
{Tăng thêm 1200 đơn vị}
{Không đổi}
{\True Giảm đi 1200 đơn vị}
{Giảm đi 2400 đơn vị}
\loigiai{
Số bị trừ giảm bao nhiêu thì hiệu giảm bấy nhiêu.
\textbf{Đáp án của câu hỏi là C}
}
\end{ex}

\begin{ex} %[3D2V64]
Một xã có $78\,500$ nhân khẩu. Sau một năm có $2400$ người chuyển đi nơi khác. Số dân của xã còn lại là
\choice
{$75\,100$ người}
{\True $76\,100$ người}
{$76\,000$ người}
{$77\,100$ người}
\loigiai{
$78\,500 - 2400 = 76\,100$ người.
\textbf{Đáp án của câu hỏi là B}
}
\end{ex}

\begin{ex} %[3D2V64]
Tính giá trị của biểu thức: $100\,000 - (45\,000 + 25\,000)$.
\choice
{$20\,000$}
{$40\,000$}
{$25\,000$}
{\True $30\,000$}
\loigiai{
Trong ngoặc: $45\,000 + 25\,000 = 70\,000$. Phép trừ: $100\,000 - 70\,000 = 30\,000$.
\textbf{Đáp án của câu hỏi là D}
}
\end{ex}

\begin{ex} %[3D2V64]
Một đại lý xăng dầu nhập về $90\,000\ell$ dầu. Tuần 1 bán $32\,000\ell$, tuần 2 bán $35\,000\ell$. Sau hai tuần đại lý còn lại bao nhiêu lít dầu?
\choice
{\True $23\,000\ell$}
{$22\,000\ell$}
{$25\,000\ell$}
{$24\,000\ell$}
\loigiai{
Tổng bán: $32\,000 + 35\,000 = 67\,000\ell$. Còn lại: $90\,000 - 67\,000 = 23\,000\text{ (\ell)}$.
\textbf{Đáp án của câu hỏi là A}
}
\end{ex}

\begin{ex} %[3D2V64]
Tìm số tự nhiên biết rằng lấy số lớn nhất có năm chữ số trừ đi số đó thì được 45 678.
\choice
{$54\,312$}
{$53\,321$}
{\True $54\,321$}
{$54\,231$}
\loigiai{
Số cần tìm: $99\,999 - 45\,678 = 54\,321$.
\textbf{Đáp án của câu hỏi là C}
}
\end{ex}

\textbf{Phần 2. Trắc nghiệm đúng sai}
\setcounterpageref{ex}{0}

\begin{ex} %[3D2V64]
Xét tính đúng sai của các phép trừ trong phạm vi 100 000:
\choiceTF
{\True Thực hiện phép trừ lần lượt từ phải sang trái.}
{$100\,000 - 1 = 9999$.}
{\True Khi cùng bớt ở cả số bị trừ và số trừ một số như nhau thì hiệu không thay đổi.}
{$85\,000 - 45\,000 = 30\,000$.}
\loigiai{
\begin{itemize}
    \item Ý a) Đúng: Quy tắc trừ.
    \item Ý b) Sai: Sai vì $100\,000 - 1 = 99\,999$.
    \item Ý c) Đúng: Tính chất của hiệu.
    \item Ý d) Sai: Sai vì $85\,000 - 45\,000 = 40\,000$.
\end{itemize}
\begin{itemize}
    \item a) \textbf{Đúng}.
    \item b) \textbf{Sai}.
    \item c) \textbf{Đúng}.
    \item d) \textbf{Sai}.
\end{itemize}
}
\end{ex}

\begin{ex} %[3D2V64]
Một bể nước chứa $80\,000\ell$ nước, người ta xả nước hai lần:
\choiceTF
{\True Lần 1 xả $25\,000\ell$ thì trong bể còn lại $55\,000\ell$ nước.}
{\True Lần 2 xả $35\,000\ell$ thì trong bể còn lại $20\,000\ell$ nước.}
{Tổng lượng nước xả ra sau hai lần là $50\,000\ell$.}
{\True Lượng nước còn lại bằng một phần tư lượng nước ban đầu.}
\loigiai{
\begin{itemize}
    \item Ý a) Đúng: $80\,000 - 25\,000 = 55\,000\ell$.
    \item Ý b) Đúng: $55\,000 - 35\,000 = 20\,000\ell$.
    \item Ý c) Sai: Sai vì $25\,000 + 35\,000 = 60\,000\ell$.
    \item Ý d) Đúng: $20\,000 : 80\,000 = \dfrac{1}{4}$.
\end{itemize}
\begin{itemize}
    \item a) \textbf{Đúng}.
    \item b) \textbf{Đúng}.
    \item c) \textbf{Sai}.
    \item d) \textbf{Đúng}.
\end{itemize}
}
\end{ex}

\textbf{Phần 3. Trả lời ngắn}
\setcounterpageref{ex}{0}

\begin{ex} %[3D2V64]
Tính kết quả: $94\,560 - 47\,280$.
\shortans{47280}
\loigiai{
\textbf{Đáp số:} $47280$\\
$94\,560 - 47\,280 = 47\,280$.
}
\end{ex}

\begin{ex} %[3D2V64]
Tìm số tự nhiên $y$, biết: $75\,000 - y = 28\,500$.
\shortans{46500}
\loigiai{
\textbf{Đáp số:} $46500$\\
$y = 75\,000 - 28\,500 = 46\,500$.
}
\end{ex}

\begin{ex} %[3D2V64]
Một nông trường có $64\,000\text{ kg}$ cà phê, đã bán đi $38\,500\text{ kg}$. Hỏi nông trường còn lại bao nhiêu ki-lô-gam cà phê?
\shortans{25500}
\loigiai{
\textbf{Đáp số:} $25500$\\
$64\,000 - 38\,500 = 25\,500\text{ (kg)}$.
}
\end{ex}

\begin{ex} %[3D2V64]
Tính nhanh: $86\,400 - 30\,000 - 26\,400$.
\shortans{30000}
\loigiai{
\textbf{Đáp số:} $30000$\\
$(86\,400 - 26\,400) - 30\,000 = 60\,000 - 30\,000 = 30\,000$.
}
\end{ex}

\begin{ex} %[3D2V64]
Tính hiệu: $100\,000 - 45\,800$.
\shortans{54200}
\loigiai{
\textbf{Đáp số:} $54200$\\
$100\,000 - 45\,800 = 54\,200$.
}
\end{ex}

\begin{ex} %[3D2V64]
Một đoạn đê dài $50\,000\text{ m}$, người ta đã kè đá được $32\,500\text{ m}$. Hỏi còn bao nhiêu mét đê chưa kè đá?
\shortans{17500}
\loigiai{
\textbf{Đáp số:} $17500$\\
$50\,000 - 32\,500 = 17\,500\text{ (m)}$.
}
\end{ex}
